% University of Oregon application CV, emphasizing LLM post-training and agents.
% Build from CV_AnhNH:
% pdflatex -jobname=CV-AnhNH-Oregon -output-directory=<build-dir> resume_oregon.tex
\documentclass{resume}

\usepackage[left=0.4in,top=0.3in,right=0.4in,bottom=0.3in]{geometry}
\usepackage{enumitem}
\usepackage{microtype}

\def\sectionskip{\smallskip}
\def\sectionlineskip{\smallskip}
\setlength{\parskip}{0.2em}
\hypersetup{
    pdftitle={Hoang Anh Nguyen -- University of Oregon CV},
    pdfauthor={Hoang Anh Nguyen}
}

\name{Hoang Anh Nguyen}
\address{+84 398 272 620 \\ Hanoi, Vietnam}
\address{
    \href{mailto:anhnh.4work@gmail.com}{anhnh.4work@gmail.com} \\
    \href{https://anhnh2002.github.io/}{Homepage} \\
    \href{https://github.com/anhnh2002}{GitHub}
}
\address{\textit{Published as Anh Nguyen Hoang and Nguyen Hoang Anh}}

\begin{document}

\begin{rSection}{Research Interests}

My research interests center on post-training large language models for long-horizon
agentic tasks. I aim to study how models learn to maintain relevant information, use
tools, and revise decisions from demonstrations and trustworthy feedback, while
generalizing to new tasks and retaining earlier capabilities. My background spans
continual learning, model adaptation, and software agents.

\end{rSection}

\begin{rSection}{Education}

\textbf{Hanoi University of Science and Technology} \hfill 2020 -- 2024\\
Bachelor of Science in Computer Science, GPA: 3.75/4.0
\hfill Ranked 15th of 296 students (Top 5\%)

\end{rSection}

\begin{rSection}{Selected Post-Training Project}

\looseness=-1
\textbf{Vietnamese Education LLM} \hfill AI for Real Action 2025\\
Proposed a weighted objective combining supervised fine-tuning (SFT) and direct
preference optimization (DPO); fully fine-tuned Qwen3-8B on 2K web and textbook
preference pairs for educational knowledge and natural, respectful Vietnamese
responses. Improved the competition's education benchmark by about 12 percentage
points over the base model.\par\looseness=0

\end{rSection}

% Highlight agent research alongside training-objective design and retention.
\begin{rSection}{Publications}

\begin{enumerate}[leftmargin=1.8em,label={[\arabic*]},itemsep=0.1em,topsep=0.1em]
    \item
    \textbf{CodeWiki: Evaluating AI's Ability to Generate Holistic Documentation
    for Large-Scale Codebases}\\
    \textit{\textbf{Anh Nguyen Hoang}, Minh Le-Anh, Bach Le, Nghi D. Q. Bui}\\
    Findings of ACL 2026
    \href{https://aclanthology.org/2026.findings-acl.288/}{[Paper]}
    \href{https://github.com/FSoft-AI4Code/CodeWiki}{[Code]}

    \item
    \textbf{Mutual-Pairing Data Augmentation for Few-Shot Continual Relation Extraction}\\
    \textit{\textbf{Nguyen Hoang Anh*}, Quyen Tran*, Thanh Xuan Nguyen*,
    Nguyen Thi Ngoc Diep, Linh Ngo Van, Thien Huu Nguyen, Trung Le}\\
    NAACL 2025
    \href{https://aclanthology.org/2025.naacl-long.205/}{[Paper]}
    \href{https://github.com/anhnh2002/mutual-pairing-data-augmentation}{[Code]}

    \item
    \textbf{Preserving Generalization of Language Models in Few-Shot Continual
    Relation Extraction}\\
    \textit{Quyen Tran*, Nguyen Xuan Thanh*, \textbf{Nguyen Hoang Anh*},
    Nam Le Hai, Trung Le, Linh Van Ngo, Thien Huu Nguyen}\\
    EMNLP 2024
    \href{https://aclanthology.org/2024.emnlp-main.763/}{[Paper]}
    \href{https://github.com/thanhnx12/CRE-via-MMI}{[Code]}

    \item
    \textbf{Beyond Repository Boundaries: Cross-Repository Graph Retrieval
    for Code Generation}\\
    \textit{Minh Le-Anh, Nam Le Hai, Quyen Tran, \textbf{Anh Nguyen Hoang},
    Linh Ngo Van, Bach Le, Nghi D. Q. Bui}\\
    Findings of EMNLP 2026
    \href{https://arxiv.org/abs/2609.09987}{[Paper]}

    \item
    \textbf{Enhancing Few-Shot Continual Relation Extraction via\\
    Pairwise Augmentation and Guided Sharpness-Aware Optimization}\\
    \textit{Hoang Thanh Dat*, \textbf{Nguyen Hoang Anh*}, Nam Le Hai*,
    Linh Ngo Van, Sang Dinh}\\
    Artificial Intelligence (AIJ), accepted
    \href{https://www.sciencedirect.com/science/article/pii/S0004370226001712}{[Paper]}

\end{enumerate}
\textit{* denotes equal contribution.}

\end{rSection}

\begin{rSection}{Research Experience: Continual Learning}

\textbf{Research Assistant} \hfill March 2024 -- April 2025\\
\textit{Continual Learning and Natural Language Processing}
\hfill \textit{Hanoi, Vietnam}\\
Advisors: Assistant Prof.\ Linh Ngo Van, HUST;
Associate Prof.\ Thien Huu Nguyen, University of Oregon

\begin{itemize}[leftmargin=1.5em,itemsep=0em,topsep=0.1em]
    \item Studied \textbf{few-shot continual relation extraction}: learning new
    relation types from limited labeled examples while retaining previously
    learned relations and pretrained language knowledge
    \item \textbf{Mutual-pairing data augmentation} (co-first author, NAACL 2025):
    designed augmentation and training objectives combining old and new relation
    examples into richer inputs without additional labels; reduced the accuracy drop to
    26.40\% on FewRel and 24.36\% on TACRED, with a 2.49\% improvement over
    CPL on TACRED
    \item \textbf{Preserving pretrained generalization} (co-first author, EMNLP 2024):
    reused pretrained LM heads through mutual-information maximization with the
    task classifier; improved final accuracy by about 2\% for CPL and ConPL,
    1\% for SCKD, and up to 3.75\% with LLaMA-2-7B
    \item Investigated the interaction between representation learning and
    continual adaptation; found that Sharpness-Aware Minimization reduced
    accuracy by over 6\% when used with LLM2Vec
\end{itemize}

\end{rSection}

\newpage
\vspace*{0.8em}

\begin{rSection}{Research Experience: Language-Model Agents}

\textbf{FPT Software AI Center} \hfill May 2025 -- Present\\
\textit{AI Research Resident} \hfill \textit{Hanoi, Vietnam}\\
Advisors: Associate Prof.\ Bach Le, University of Melbourne;
Assistant Prof.\ Nghi D. Q. Bui, VinUniversity

\begin{itemize}[leftmargin=1.5em,itemsep=0em,topsep=0.1em]
    \item \textbf{CodeWiki} (first author, Findings of ACL 2026): built an LLM-agent
    framework that generates repository-level documentation across seven
    programming languages through hierarchical context decomposition and
    bottom-up synthesis of module explanations
    \item \textbf{CodeWikiBench}: built hierarchical evaluation rubrics and
    agentic LLM-judge evaluation; CodeWiki achieves a 68.79\% quality score,
    +4.73 points over DeepWiki and +21.66 points over OpenDeepWiki
    \item \textbf{FALCON} (under review): refined execution-grounded checks that
    reject fabricated evidence in generated attack tests and link reported events
    to target execution; the framework improved CWEval Func-Sec@1 by 25.2 points
    as repair feedback
    \item \textbf{RustPrint} (preprint): co-developed documentation-guided C-to-Rust
    migration; with Kimi-K2, preserved 93.26\% of evaluated features across eight
    C repositories, compared with 52.52\% for a Claude Code baseline
    \item \textbf{SecPRBench} (under review, with Prof.\ Lingming Zhang, UIUC):
    curated FFmpeg, FreeType, and SQLite tasks for a benchmark of 265
    vulnerability-introducing code changes paired with their repairs
\end{itemize}

\end{rSection}

\begin{rSection}{Preprints and Under Review}

\begin{enumerate}[leftmargin=1.8em,label={[\arabic*]},itemsep=0.1em,topsep=0.1em]
    \item
    \textbf{Secure Agentic Coding through Counterexample-Grounded
    Feedback (FALCON)}\\
    \textit{Tue Le, \textbf{Anh Nguyen Hoang}, Gabriel Sasseville, Doseok Jang,
    Pierre-Luc Bacon, Nghi D. Q. Bui}\\
    Under review
    \href{https://arxiv.org/abs/2610.25916}{[Paper]}
    \href{https://github.com/FSoft-AI4Code/FALCON}{[Code]}

    \item
    \textbf{Can Agentic Code Review Catch Vulnerabilities Before They Are Merged?}\\
    \textit{Yihan Yang, Chenyuan Yang, Zijie Zhao, \textbf{Anh Nguyen Hoang},
    Ziqi Zhang, Weidong Wang, Zhiheng Lyu, Lingming Zhang}\\
    Under review
    \href{https://secprbench.github.io/paper.pdf}{[Paper]}
    \href{https://github.com/secprbench/secprbench}{[Code]}

    \item
    \textbf{RustPrint: Documentation-Guided Agentic Codebase Migration from C to Rust}\\
    \textit{Minh Le-Anh, \textbf{Anh Nguyen Hoang}, Bach Le, Nghi D. Q. Bui}\\
    Preprint, 2026
    \href{https://arxiv.org/abs/2605.14634}{[Preprint]}
\end{enumerate}

\end{rSection}

\begin{rSection}{Industry Experience}

\textbf{Oraichain Labs} \hfill May 2024 -- April 2025\\
\textit{AI Engineer} \hfill \textit{Hanoi, Vietnam}
\begin{itemize}[leftmargin=1.5em,itemsep=0em,topsep=0.1em]
    \item Led development of \textbf{Cupiee}, an emotional AI companion platform
    with 300,000+ users and about 100,000 messages per day
    \item Built its voice-agent companion; 30-day retention rose from 13\% to
    47\% following release
\end{itemize}

\end{rSection}

\begin{rSection}{Honors and Awards}

\textbf{Champion, AI for Real Action 2025} (Team Converged) \hfill January 2026\\
First prize among nearly 400 teams in Vietnam's first national AI competition
\begin{itemize}[leftmargin=1.5em,itemsep=0em,topsep=0.1em]
    \item \textbf{Techcombank Overseas PhD Study Sponsorship}: up to USD 333,000
    for overseas PhD study, awarded following the competition (April 2026)
    \href{https://anhnh2002.github.io/Letter_of_Offer_N.H.Anh.pdf}{[Offer Letter]}
\end{itemize}

\end{rSection}

\begin{rSection}{Selected Open-Source Software}

\textbf{\href{https://github.com/anhnh2002/XTTSv2-Finetuning-for-New-Languages}
{XTTSv2-Finetuning-for-New-Languages}} \hfill 200+ GitHub stars\\
Toolkit for adapting XTTSv2 to new languages, including a released Vietnamese model;
sliding-window streaming reduces time to first audio from about 1 s to 200 ms

\end{rSection}

\begin{rSection}{Languages}
Vietnamese (native), English (IELTS 7.0, TOEFL 5.0)
\end{rSection}

\end{document}
